\section*{Chapter \ref{ch:iv:applications} --- Applications}

\begin{solution}{iq:iv:applications:1}
Say what was tested, on what, with what result against what baseline, and what failed: ``Tested whether
one-hour order-flow imbalance on two exchanges predicts the next hour's return in 20 crypto pairs, 2022--2025;
walk-forward hit rate 51.2\% against 50.0\% (t = 1.4), not significant after costs; Python.'' The rewrite
gives the interviewer three things to ask about (the data, the validation, the costs) and shows the candidate
reports a null result honestly.
\iqlookfor{a verb, an object, a number with its baseline, and honesty about the outcome.}
\end{solution}

\begin{solution}{iq:iv:applications:2}
Four sentences, in the order of relevance to the role: what you are now (degree, year, subject); the one
piece of evidence closest to trading (an internship, a competition, a project with numbers); what that taught
you about decisions under uncertainty; why this role now. Leave out the chronology, every item the
interviewer can read, and anything you cannot discuss for five minutes. Stop at sixty seconds and let the
interviewer choose what to open.
\iqlookfor{selection and order by relevance, brevity, and a clear reason for the application.}
\end{solution}

\begin{solution}{iq:iv:applications:3}
The CV should have listed only languages the candidate can be interviewed in for thirty minutes, with the
others under ``also used'' or not at all. Now: say plainly that Haskell was a short course project, that the
strong languages are C++ and Python, and offer to write the solution in either while explaining how it would
look in a functional style (recursion instead of loops, immutable data, a fold). Pretending costs more than the
admission; the interviewer chose Haskell to find out.
\iqlookfor{honest calibration of one's own skills, and a constructive alternative.}
\end{solution}

\begin{solution}{iq:iv:applications:4}
With 400 equal teams and noise 0.01, the expected best public score exceeds the true score by
$0.01 \times m_{400} \approx 0.01 \times 2.97 \approx 0.030$ (\cref{prop:iv:applications:max}). A drop from third
to fortieth is what noise alone produces when the public sample is small; the private leaderboard, scored on
more data, is the better measurement. In the interview: say so, say what you did to avoid fitting the public
board (cross-validation that mimicked the private split, a limit on submissions used for selection), and what
you would do differently.
\iqlookfor{the expected maximum of noisy scores, and a mature reading of one's own result.}
\end{solution}

\begin{solution}{iq:iv:applications:5}
The fee is $0.20 \times 150\,000 = 30\,000$. At 160\,000 it would be 32\,000, a gain of 2\,000, but with a
90\% chance of closing the expected fee is $0.9 \times 32\,000 = 28\,800$, which is 1\,200 below the certain
30\,000. On these numbers the recruiter prefers that you do not ask. The advice is not dishonest for that; it
is predictable. Ask the recruiter for facts (the band, what the firm has paid similar hires), not for a
recommendation, and negotiate the number yourself (\cref{ch:iv:offers-compensation-and-non-competes}).
\iqlookfor{an expected-value reading of an intermediary's incentives.}
\end{solution}

\begin{solution}{iq:iv:applications:6}
Make the CV machine-readable: plain text in a standard layout, conventional section names, the role's
vocabulary where it is true, dates in one format, no content only in images or tables. Answer form questions
precisely. Do not stuff keywords or hide text: automated screens are followed by people, who reject it. On
rights: in New York City such tools must have had a bias audit and candidates must be notified; in the EU
recruitment tools are high-risk systems under the AI Act, with requirements applying from December 2027
(\cref{dat:iv:applications:law}); elsewhere, ask the firm what the tool does and whether an alternative
process exists, which some firms offer as an accommodation.
\iqlookfor{practical legibility, no gaming, and awareness that candidates have rights that vary by place.}
\end{solution}

\begin{solution}{iq:iv:applications:7}
Over five years an annualised Sharpe ratio estimate has standard deviation about $1/\sqrt5 \approx 0.447$, so
the best of sixty worthless variants is expected near $2.32 \times 0.447 \approx 1.04$. The reported 2.4 is
about $(2.4 - 1.04)/0.447 \approx 3.0$ standard deviations of one estimate above that, which is suggestive but
not conclusive, since the variants are correlated and the costs may be understated. Defend it by saying what
the sixty variants were, how the final one was chosen, what the out-of-sample or later-period result was, and
the deflated Sharpe ratio (One Quant Book~4, chapter~12). If none of that exists, retract the number and keep
the line as a description of the study.
\iqlookfor{the multiple-testing correction, done in the head, and intellectual honesty about one's own CV.}
\end{solution}

\begin{solution}{iq:iv:applications:8}
Ask her what the team and the role are like, what the firm values that its website does not say, and whether
she knows your work well enough to refer you; if she does, ask her to submit or endorse the application and
to say what she actually knows about you. Do not ask for interview questions: it puts her in breach of her
employer's trust and the questions will change anyway. A referral usually moves the application past the
first screen or to a person's desk, and it raises the prior the screen starts from; it does not change the
later stages, which are the same for everyone (Burks et al., 2015, on why firms value referrals).
\iqlookfor{using a referral for information and endorsement without compromising the referrer.}
\end{solution}
